% =============================================================================
%  Conclusion générale
%
%  Aucun chiffre n'est introduit ici qui ne figure déjà, avec son contexte et
%  sa marge d'erreur, au chapitre 5. La conclusion synthétise ; elle ne
%  renforce pas.
% =============================================================================

\chapter*{Conclusion générale}
\addcontentsline{toc}{chapter}{Conclusion générale}
\label{chap:conclusion}
\markboth{CONCLUSION GÉNÉRALE}{CONCLUSION GÉNÉRALE}

% =============================================================================
\section*{Ce qui était demandé, ce qui a été livré}
% =============================================================================

L'énoncé d'EURAFRIC Information demandait un système d'optimisation de
portefeuille piloté par apprentissage automatique, capable d'apprendre les
corrélations dynamiques entre actifs, opérant sous des contraintes réalistes de
gestion, et évalué dans un cadre rigoureux de \emph{backtesting} hors
échantillon.

Le livrable est une chaîne complète et reproductible~: collecte et validation de
données provenant de cinq sources hétérogènes~; une couche de préparation en
trois niveaux, validée à l'écriture~; un moteur de simulation temporelle dont
l'absence de vision du futur est vérifiée par des tests plutôt qu'affirmée~;
quatre familles de modèles ajoutées une à une, du rétrécissement de covariance à
la détection de régime, de la prédiction de rendement à l'optimisation sensible
aux coûts~; un protocole d'évaluation comportant une sélection strictement
antérieure, un segment de test gelé, des \textbf{comparaisons pairées} et une
correction pour le nombre d'essais~; trois cartes de modèle, une couche
d'explicabilité et une instrumentation des replis~; et deux surfaces de
consultation — une application web à deux pages et une interface REST en lecture
seule — partageant une couche de données unique. L'ensemble est couvert par
\textbf{plus de 780~tests} exécutés hors ligne et par un contrôle du PDF rendu.

% =============================================================================
\section*{Ce que le projet établit}
% =============================================================================

Les résultats sont classés selon leur force probante, et cette distinction est
elle-même un résultat de la démarche.

\paragraph{Exploré.} Le détecteur de régime, non supervisé et n'ayant jamais
reçu de liste de crises, est plus souvent dans son état étiqueté « baissier » à
l'intérieur des cinq fenêtres de crise étudiées (91,7~\%) qu'en dehors (29,2~\%),
soit un rapport de \textbf{3,13}. Le diagnostic par test des signes ($p = 0{,}031$)
est conservé pour transparence, mais ne constitue pas une preuve confirmatoire~:
l'échantillon ne comporte que cinq événements et l'état « baissier » est défini
par son rendement moyen plus faible. La portée défendable est plus étroite~: le
modèle produit une lecture causale de l'état de marché associée à ces fenêtres ;
il ne démontre pas qu'agir sur cette lecture améliore la performance économique.

\paragraph{Observé, non démontré.} Sur le portefeuille visé, libellé en dirhams
au taux de référence officiel de Bank Al-Maghrib et non couvert contre le risque
de change, l'écart ponctuel entre la stratégie à régimes et la meilleure approche
classique est de \textbf{$-10{,}47$~\%} de ratio de Sharpe net~; il est également
défavorable sur l'univers ETF, à $-1{,}6$~\%. Ces écarts ne démontrent pas une
supériorité de l'approche classique~: aucun test pairé de cette différence n'est
présenté. Le classement est par ailleurs sensible au protocole d'évaluation et à
la fenêtre hors échantillon associée.

\paragraph{Testé, non établi.} La réserve précédente n'est plus une réserve~: le
test manquant a été mené. Sous sélection strictement antérieure — l'entraînement
ne pouvant plus consulter des dates postérieures à sa fenêtre de validation — les
\textbf{huit} comparaisons pairées effectuées sur la fenêtre de test gelée ont un
intervalle contenant zéro ($p$ compris entre 0,066 et 0,939). \textbf{Aucune
surperformance n'est établie}, ni pour les modèles de signal contre la stratégie à
régimes, ni contre l'équipondéré, sur aucun des deux univers.

Un cas résume l'intérêt de la démarche. Sur \texttt{etf\_2017}, le signal RF
calibré obtient 1,102 contre 1,093 pour la stratégie à régimes~: en estimation
ponctuelle il « bat » la référence. Le test pairé donne pourtant $p = 0{,}326$.
Le classement ponctuel, celui qu'on aurait rapporté sans ce test, n'est pas
soutenu par les données.

Ce n'est pas pour autant une preuve d'équivalence~: ne pas rejeter n'est pas
accepter l'hypothèse nulle. L'établir demanderait un test d'équivalence contre une
marge fixée à l'avance, qui n'a pas été mené. Ces réserves sont affichées dans le
rapport, le tableau de bord et l'interface de programmation.

\paragraph{Observé sur cinq fenêtres.} Pendant les cinq crises, l'optimisation sous
contrainte a systématiquement mieux protégé le capital que la diversification
naïve, avec un délai de retour au sommet environ deux fois plus court.
L'attribution est faite honnêtement~: ce gain revient à la contrainte de position
et au modèle de covariance, non à la couche de régime, qui dans trois de ces cinq
fenêtres produit exactement la même allocation qu'une stratégie classique.

\paragraph{Découvert.} Le plafond de 25~\% par actif, retenu comme règle métier et
non comme dispositif de modélisation, s'est révélé le régulariseur le plus
efficace du système~: le relâcher coûte 10,1~\% de ratio de Sharpe, davantage que
l'écart entre deux modèles quelconques testés sur cet univers. Ce résultat est la
reproduction, sur nos propres données, d'un théorème connu — une contrainte de
position active équivaut mathématiquement à un rétrécissement de la matrice de
covariance.

\paragraph{Réfuté.} Cinq pistes indépendantes visant à dépasser la ligne de base
ont été explorées et fermées~: la prédiction de rendement calibrée honnêtement,
un historique cinq fois plus long, l'ajout de données fondamentales, la
re-sélection à chaque repli temporel, et le conditionnement de la contrainte au
régime. Le constat qui émerge de leur convergence est plus intéressant que chaque
échec pris isolément~: \textbf{aucune de ces pistes n'établit un avantage
incrémental de portefeuille dans les protocoles évalués}. Deux de ces études
montrent en outre, séparément, que \emph{précision de prédiction et
performance de portefeuille sont deux objectifs distincts}~: le coefficient
d'information a doublé sans que le ratio de Sharpe suive, et il a même reculé
dans un cas.

% =============================================================================
\section*{Ce que le projet apporte au-delà des chiffres}
% =============================================================================

La contribution la plus durable de ce travail n'est probablement pas une
performance, mais un \emph{dispositif}.

Deux incidents l'illustrent mieux qu'un discours. La découverte que les
dividendes marocains manquaient a réduit de plus de moitié le gain historique.
Puis la correction du numéraire, au taux officiel Bank Al-Maghrib, a renversé
son signe. Les chiffres corrigés ont remplacé les anciens sur chaque surface et
les affirmations antérieures ont été retirées. Ces corrections n'ont été
possibles que parce que la chaîne était reconstructible de bout en bout ; elles
auraient probablement été masquées dans un projet où les résultats sont produits
à la main.

De là vient le principe qui structure l'ingénierie du projet~: \textbf{une règle
importante n'est pas documentée, elle est rendue exécutable}. L'impossibilité de
consulter le futur, l'interdiction d'écrire un chiffre de performance en dur dans
une interface, la cohérence entre documents publiés, la traçabilité de chaque
fonction vers le problème qu'elle traite, l'hermétisme réseau de la suite de
tests~: chacune de ces règles est un test qui échoue si on l'enfreint. Les trois
incidents les plus coûteux du projet ont tous consisté en une règle respectée par
discipline jusqu'au jour où elle ne l'a plus été, sans que rien ne le signale.
Chacun a donné lieu à un garde-fou automatique, validé en rejouant l'incident
réel. Les cartes de modèle et les traces d'explication complètent ce dispositif :
la stratégie de référence est expliquée par ses décisions, les challengers par
des attributions additives vérifiées, et le monitoring est explicitement prêt
mais non actif tant que ses préconditions de publication atomique et de rollback
ne sont pas satisfaites.

% =============================================================================
\section*{Limites}
% =============================================================================

Elles sont énoncées ici plutôt que laissées à découvrir.

\begin{itemize}
  \item \textbf{La taille et la composition de l'échantillon limitent la
        puissance, non toute conclusion.} \texttt{full\_2021} couvre environ
        cinq années~: cela ne permet pas de départager solidement de petits écarts
        entre stratégies. Le projet conclut toutefois, sur les protocoles
        évalués, qu'aucun avantage ML robuste n'est établi.
  \item \textbf{L'univers BVC canonique commence en 2021.} Un historique marocain
        profond 2005--2024 a bien été testé séparément~: il a amélioré la qualité
        de prédiction sans établir de gain de portefeuille. Il n'est pas intégré
        à la release parce que source, raccordement récent, dividendes et
        opérations sur titres doivent encore être contractualisés et reproduits.
        L'univers ETF, lui, remonte à novembre 2004 et couvre 2008, 2020 et 2022,
        sans se substituer à un historique BVC canonique long.
  \item \textbf{Un rendement de dividende sur quatre a dû être estimé}, la source
        publique ne le publiant pas. L'analyse de sensibilité correspondante est
        documentée, mais la donnée reste une estimation.
  \item \textbf{La correction pour le nombre d'essais est établie} sur les 240
        configurations atteignables de la recherche contre la référence primaire.
        Huit comparaisons externes demeurent non corrigées à ce niveau, et les
        résultats contre l'équipondéré restent exploratoires~: cette frontière
        borne explicitement les affirmations autorisées.
  \item \textbf{Le système n'a pas été exposé à des conditions d'exploitation
        réelles}~: ni exécution d'ordres, ni flux temps réel, ni couverture du
        risque de change dirham/dollar. Le monitoring est instrumenté hors ligne
        mais reste volontairement non actif derrière ses portes de release.
\end{itemize}

% =============================================================================
\section*{Perspectives}
% =============================================================================

Quatre suites se dégagent, classées par rapport valeur/effort.

\textbf{Industrialiser l'historique marocain profond.} L'expérience a déjà écarté
l'idée qu'ajouter du seul historique de prix promette mécaniquement un avantage
de portefeuille. Elle reste néanmoins précieuse pour la puissance statistique,
la couverture de 2008 et la représentativité de l'univers BVC. La suite est un
travail de gouvernance et de validation des données~: contrat de source,
rapprochement sur le recouvrement, dividendes et opérations sur titres, puis
réévaluation complète. Ce n'est pas une promesse de performance.

\textbf{Faire valider le modèle indépendamment et activer le monitoring.} Une
validation extérieure doit challenger le protocole, les hypothèses de coûts et
les artefacts de release. Ensuite seulement, le monitoring pourra être activé
après publication atomique, vérification de manifeste, verrouillage des releases
et exercice de rollback.

\textbf{Calibrer la pénalité de rotation par modèle.} Les mesures ont établi
qu'une valeur globale unique dégrade le bon modèle tout en corrigeant le mauvais.
La rendre spécifique à chaque modèle est un travail de calibration circonscrit,
au protocole déjà disponible.

\textbf{Adoucir le basculement de régime.} Le basculement franc actuel produit un
pic de transactions à chaque frontière de régime. Un mélange continu, pondéré par
la probabilité de régime, est une alternative documentée dès la conception et non
encore mesurée.

% =============================================================================
\section*{Bilan}
% =============================================================================

Ce projet aura moins appris à l'équipe comment faire gagner un modèle qu'à
\emph{savoir si un modèle gagne}. C'est un apprentissage inconfortable~: il
consiste, l'essentiel du temps, à construire les instruments qui feront tomber
les résultats auxquels on tenait — un test qui interdit de voir le futur, un
intervalle de confiance qui absorbe l'écart qu'on voulait annoncer, une donnée
corrigée qui divise un gain par deux à trois semaines de la remise.

C'est pourtant, dans un domaine où les performances publiées sont notoirement
difficiles à reproduire, la seule compétence qui se transporte d'un sujet à
l'autre. Le prototype livré à EURAFRIC Information ne prétend pas battre les
méthodes classiques, ni recommander ou exécuter une allocation~: il expose les
écarts historiques observés, leurs limites et les conditions de leur
reproduction. Il constitue un socle de recherche et d'analyse, pas un outil de
gestion en production.
